% !TeX root = ../papers/nc-p3.tex
\section{三次縮約の三支持からの部分代数と低い反対階数の全次数判定}
\label{sec:ternary-point-subalgebras}
有限Frobenius窓の標準四テンソルを $w$ とし、全辺の左右の射影ファイバーが空または一点とする。
点グラフの次元は問わない。一文字を点 $a$ で縮約した三テンソルを巡回順に読む。
\begin{theorem}\label{thm:ternary-point-subalgebras}
ある縮約 $u=w(a)$ の三つの一因子支持が全て三次元であるとする。
各因子の支持に制限した $3\times3\times3$ テンソル $u$ について、次は同値である。
\begin{enumerate}
\item 元の四窓が標準Koszul左右AS正則である。
\item 三方向のいずれの非零一文字縮約にも階数1の行列がない。
\end{enumerate}
従って一つの縮約の三つの射影平面で二次小行列式の共通零点を調べれば足りる。
このとき $A$ は三生成二次AS正則部分代数の上で左右自由であり、
判定に失敗すると指数成長する。
\end{theorem}
\begin{proof}[三文字の列の構成]
最初の三支持を順に $H_i,H_{i+1},H_{i+2}$ とする。添字を四移動で読めば
$H_i=R_{i-1}(a)$、$H_{i+2}=R_{i+2}(a)$ であり、
その直交点 $p_i,p_{i+2}$ は $a$ の次点、前点である。
$u$ の末尾と先頭の縮約像はそれぞれ $U_i\subset H_i\otimes H_{i+1}$,
$U_{i+1}\subset H_{i+1}\otimes H_{i+2}$ という三次元の二次関係空間である。
各 $U_j$ は $R_j$ のうち先頭が $H_j$ に支持される部分および末尾が $H_{j+1}$ に
支持される部分に等しい。実際それぞれの評価核に含まれ、その核の次元はファイバー仮定で高々3である。
特に中支持の直交点 $p_{i+1}$ も共有点グラフに属する。

$p_{i+1}$ の唯一の前点 $b$ で $w$ を縮約すると、位置 $i+1$ から始まる三テンソルの
先頭は $H_{i+1}$ に支持され、末尾の支持も三次元である。その末尾を縮約した像は上の
$U_{i+1}$ と等しいので、中央の支持はちょうど $H_{i+2}$ である。
こうして次の $H_{i+3}$ が定まり、逆向きも同様に延長できる。
得られた全列で $U_j$ は上の左右の評価核であり、三次重なりは既知の三テンソルの直線である。
後者は、元の四文字代数の三次重なりのうち $p_j$ での先頭縮約が零になる部分が一次元であるためである。
実際この部分は $w$ の末尾縮約のうち $p_j$ の前点を使うものだけである。

抽象二次代数 $P=T((H_j))/(U_j)$ の二次双対は従って $(1,3,3,1)$ の有限Frobenius型である。
三次テンソルの各一因子縮約は相補関係への同型、四次重なりは非零なら三次テンソルの片端支持が
一次元になってしまうので零である。次の三テンソルは前のものの巡回回転と、余った一因子の
基底変更で一致する。実際、両者は共通の $U_{i+1}$ への同型を与える。
従って全三位置の縮約階数は最初の $u$ で検査できる。
\end{proof}

\begin{proof}[自由性と有限判定]
各始点の $p_j$ における右点加群を $M_j$ とする。
命題\ref{prop:commuting-three-point-graph}の切断計算と同じく
$M_j=e_jA/H_jA$ であり、$H_j\otimes e_{j+1}A\to e_jA$ の関係は
次数2の高々三個で生成される。この計算に三文字の可換性は必要なく、
元の候補複体上の点列での階数 $1,3,3,1$ を用いる。
二次には独立な $U_j$ が既にあるからそれらが生成する。
$P$ 単純の自由表示の最初の二写像と比較して $\operatorname{Tor}^P_1(k_j,A)=0$。
各終点の有界下加群の極小表示から $A$ は左 $P$ 自由で、各長さに一個の生成元がある。
逆向きの点加群によって右自由性も従う。特に $P$ は $A$ に埋まる。

$P$ の三窓について、補題\ref{lem:ternary-point-graph-growth}と
補題\ref{lem:ternary-finite-graph-test}の証明を添字付きで適用する。
具体的には、三テンソルを順に回転する表示では、一つの点での左評価と右評価は
同じ二文字行列の転置または定数基底変更である。階数1なら、その両端像を割った二文字には
二次元ずつが残り、関係は一本が消えて二本以下である。
三位置を一期として独立に繰り返し、長さ $3m$ で $2^m$ 以上となる。
逆に全て階数2以上なら三位置の両側射影は同じ行列式軌跡をスキームとして共有し、
全て平面または平面三次曲線の同型のグラフになる。平面なら可換多項式型である。
三次の場合はグラフに沿う束 $L_j$ を順次使えば、上記補題と同じ分裂局所完全な
階数 $1,2,1$ の評価複体になる。正の偏極積の $H^1$ は消え、長さ3の左端の
$H^1(\mathcal O)$ だけが残る。よって全切断環を与える商への核は各位置で
左右とも正則な三次の一元で生成され、$P$ のHilbert関数は三変数多項式環のものとなる。
左右のsocle消滅、Euler式および三次対合から $P$ は標準三次元Koszul左右AS正則となる。
これらの計算に同じ束の反復や普通の次数付き表示を仮定しない。

最後に $P$ が正則なら左右自由性と各長さ一個の自由生成元より $H_A=(1-z)^{-4}$。
任意の位置の $H_j$ の二次元部分空間には右乗法が単射の文字が存在する。
平面の場合は明らかであり、三次曲線の場合は独立な二一次式が共に恒等零となる既約曲線成分はないので、
全ての随伴点を避ける一次式を取れる（平面三次は純である）。
その右乗法は切断商で単射であり、先の正則三次元を左側にくくり出す次数降下により $P$ でも単射である。
右 $P$ 自由性から $A$ でも同じである。従って点加群の任意の尾部から $A$ への右Homは零である。
実際、各位置の点加群の元は $H_j$ の二次元以上で消える。
$e_jA_+$ の右線型写像を $e_jP_+$ に制限すれば、三次元AS性と自由性で
左乗法に延びる。それを引いた写像は $H_jA$ を消して点加群の尾部を通るので零である。
これで右単純から $A$ への次数0,1のExtが消え、左側も同様である。
四文字候補複体のFrobenius双対およびHilbert関数のEuler式で正則性が従う。
\end{proof}

\begin{corollary}\label{cor:low-opposite-point-graph-as}
点グラフの四辺の左右射影ファイバーが空または一点なら、反対階数 $\rho\le7$ の
全候補で標準AS正則性が有限に決まる。$\rho<6$ は全て除外される。
$\rho=6,7$ の場合は三支持が三次元の縮約の有無を調べ、なければ除外、
あれば一つの縮約で前定理を適用する。
特に四辺が有限な枝は $\rho\le7$ の全域で、全次数の正則性と印を指定しない添字保持同型・周期が
有限判定となる。
\end{corollary}
\begin{proof}
六双一次式の零点は空でない。実際、Segre埋込みの六次元射影多様体と余次元6の線型空間は交わる。
標準なら任意の零点で三次縮約の両端支持は3であり、中支持は少なくとも3である。
中支持2は補題\ref{lem:middle-binary-pencil}の両端次元の不一致に反し、
中支持1なら同じ中央文字での連続する純粋関係ができる。
反対切断の $\rho$ 次元を経る二つの評価行列は、辺の点対で積が零である。
従って両点での中支持の和は $\rho$ 以下となる。
$\rho\le7$ なら少なくとも一方が3である。有限点の場合の同型と周期には
定理\ref{thm:finite-all-window-keys}を適用する。
\end{proof}

\begin{theorem}\label{thm:finite-frobenius-low-opposite-as}
四辺が有限な標準有限Frobenius候補では、反対階数によらず全ての $R_i$ が
非零の可分解二テンソルを含まない。この枝で $\rho<6$ の候補は存在せず、
$\rho=6,7$ なら全ての候補が標準Koszul左右AS正則である。
\end{theorem}
\begin{proof}
系\ref{cor:finite-edge-point-sequences}の共通長さ20のスキームを $X$ とし、
直線束の自明化後の四文字の切断空間を $V_i\simeq U_i\subset D=H^0(\mathcal O_X)$ と書く。
定理\ref{thm:finite-square-chart}の六双一次式の完全交叉とtraceの議論は、
この共通グラフが成立してから先には元代数のAS性を使わない。
従って各 $\operatorname{Sym}^3 U_i\to D$ は同型、$R_i$ は積の核、
$U_{i+2}^2$ は $U_i^2$ の単元倍である。

もし非零の $a\in U_i,b\in U_{i+1}$ について $ab=0$ なら、
$a U_i^2$ と $b U_{i+1}^2$ は三次の単射性により各十次元である。
$(a)\subset\operatorname{ann}(b)$ かつ $\dim\operatorname{ann}(b)=20-\dim(b)$ なので、
両主イデアルが次元10となる。$Y=\operatorname{Spec}(D/(a))$ を三つ目の四文字 $U_{i+2}$ で見る。
この空間の二次切断の制限像は次元6以下である。実際 $U_i^2$ の制限核は
四次元の $a U_i$ を含み、単元倍で核の次元は変わらない。三次の制限像は全長10である。

この $Y$ は平面に入る。入らないとすると斉次座標環のHilbert関数は
$1,4,h_2,10$ で $h_2\le6$ となる。全ての点で非零の一次式を取って割ると、
非零因子なので差分のHilbert関数を得る。次数2が $h_2-4\le2$ で次数3が $10-h_2\ge4$ となるのは不可能である。
実際、初項イデアルの標準単項式は約数で閉じており、二次の標準単項式高々二つからは
三次のものは高々二つしか作れない。従ってある $0\ne c\in U_{i+2}$ が $Y$ 上で消え、
$c\in(a)$、ゆえに $bc=0$ となる。すると $a\otimes b\otimes c$ が三次重なりに入る。
有限Frobenius対合によりこれは標準四テンソルのある非零一文字縮約である。
その両端支持が一次元ということは、その文字の点で左右の関係評価が階数1ということであり、
有限ファイバーに反する。これで可分解関係は全て排除された。

辺の各点で三次縮約の両端支持は3である。もし中支持が2なら、
中央因子で縮約する $3\times3$ 行列の鉛筆はある非零パラメータで特異になる。
その片側核で元の三テンソルを縮約すれば、残る中央支持が一次元の非零二次関係を与え、
先に得た結論に反する。中支持1も同様に不可能である。
従って上の反対階数の和の不等式はAS性を仮定せず適用でき、$\rho\le7$ なら三支持が3の元が存在する。
その三方向の縮約のうち両端からのものは既に $R_j$ 内にある。
中央からの縮約も階数1になれないことを確認する。もしそうなら、
構成した三文字部分代数の第三の関係空間はその巡回回転を含む。
この空間も延長した $U_{j+2}\subset R_{j+2}$ に可逆基底変更で等しいから矛盾する。
ゆえに前定理により正則である。
\end{proof}
